\section{A retention protocol that survives a publishing race}
\label{exec17:i4-pins}
The preceding reachability rule states which artifacts are needed at one instant. A collector must additionally preserve that property while a report becomes visible. Consider a generated deposit report whose immutable dependency descriptor points to the input balance extract. Before publication, neither object need be reachable from an accepted report. A publisher can be preparing a legitimate result precisely when a collector takes its snapshot. The substantive operating question is whether the proposed coordination rule prevents a published report from losing its input, without assuming that every process which started will eventually finish.

For a finite model, let $V$ be the set of present immutable object identities, $R$ the accepted report roots, and $P$ the roots pinned by admitted but unfinished publishers. Each descriptor has an immutable, complete list of dependencies. Write $\operatorname{cl}(S)$ for $S$ together with all objects reachable through these lists. The operating invariant is
\begin{equation}
 \operatorname{cl}(R\cup P)\subseteq V.
 \label{exec17:i4-live-invariant}
\end{equation}
This is a storage-availability condition for declared dependencies. It cannot verify that a descriptor lists all the scientific inputs actually used, that a balance extract is accurate, or that a customer record may be retained. Those questions require separate provenance and custody evidence. The following protocol is a proposed local design and finite reference, not an implemented collector or a theorem about arbitrary distributed storage~\citep{exec17_retention_reference}.

\subsection{Pin acquisition must agree with the deletion decision}
An admitted publisher needs a reservation that protects a complete dependency closure. Under the proposed local rule, pin acquisition checks that all required objects exist and adds the pin while holding the same metadata coordination boundary used by deletion. The publisher subsequently installs the report root before releasing its pin. Publication, pin release and the collector's final deletion decision use that boundary too. An object being staged outside the store is not yet promised to any consumer. If collection wins before admission, the publisher must reconstruct or reacquire the missing objects and seek admission again; it cannot claim a valid pin on already missing bytes.

Let $M$ be the closure marked at the collector's earlier snapshot. Immediately before deletion, re-read the current roots and pins under coordination and remove only
\begin{equation}
 D=V\setminus\bigl[M\cup\operatorname{cl}(R\cup P)\bigr].
 \label{exec17:i4-delete-set}
\end{equation}
An old mark can keep extra objects, which costs storage, but the current closure prevents a newly protected object from being collected. The final dependency check and deletion must be indivisible with respect to pin acquisition and root publication in this model. Checking first, releasing the lock, and deleting later reintroduces the same race. A scalable implementation may use another mechanism, but it owes an equivalent argument; naming a snapshot or a quarantine does not establish that argument.

The invariant has a short inductive proof. Begin with all current roots and pins present. A successful pin acquisition checks its entire closure, so adding that pin preserves the inclusion. Rejected admission changes neither roots nor pins. Publication adds a root whose closure was protected by a pin, and occurs before that pin is removed. Releasing a pin can only reduce the left side of the inclusion. A snapshot changes no live object. Finally, deletion by the displayed rule removes nothing in the current closure. Every permitted transition therefore preserves the invariant. This proof depends on complete immutable descriptors and the shared coordination boundary; it does not follow merely from individually atomic file renames.

\subsection{Ten schedules separate safe rejection from broken publication}
The reference has two objects: a report and its input, with one edge from the report to the input. A publisher performs pin, publish and unpin in that order. A collector performs snapshot then sweep. There are $\binom{5}{3}=10$ order-preserving interleavings. This denominator follows by choosing the three publisher positions; it does not assume equally likely runtime schedules. The reference enumerates all ten and retains the state after each event~\citep{exec17_retention_reference}.

Three interleavings violate the invariant when sweep uses only its stale snapshot. One witness is snapshot, pin, publish, sweep, unpin. The snapshot sees no roots. Pin acquisition then sees both objects, and publication exposes the report. A stale sweep deletes both objects despite that successful publication. A second order, snapshot, pin, sweep, publish, unpin, breaks the pin before publication. In the third, snapshot, pin, publish, unpin, sweep, the published root still fails to protect its objects because the collector never rechecks it. The defect is independent of whether the numerical report was computed correctly.

With the current-closure recheck, all ten schedules preserve the invariant. Nine publish successfully. In the remaining order, snapshot and sweep both occur before pin acquisition; both unprotected objects have been collected, so admission rejects and no report becomes public. This is an availability failure for that attempt, yet it is the correct integrity outcome. Treating a rejected attempt as an unsafe publication would confound two distinct criteria. Conversely, counting only the nine successful attempts and suppressing the rejection would overstate availability. A subsequent implementation trial should retain both denominators.

The example does not estimate a crash probability or report a measured three-in-ten failure rate. It is exhaustive only for this five-event model. Real implementations add content writes, descriptor writes, error returns, multiple publishers and filesystem failures. The proof identifies which transitions require a shared boundary and which additional events need refinement. The next low-cost experiment can therefore test a small disposable store with deliberate pauses at these boundaries, rather than run an expensive bank-data study to rediscover a deterministic race.

\subsection{A crashed publisher creates a liveness obligation}
Pins introduce a second question: who releases a pin after its publisher disappears? Leaving a stale pin in place protects its complete closure indefinitely. The reference's two-object example retains both objects, even if no accepted report will ever need them. This is safe under the inclusion invariant but can exhaust storage. Safety and eventual reclamation are separate requirements; a system should not claim both from the same reachability calculation.

Deleting a pin after a timeout is insufficient if a paused publisher can later resume. Once the expired pin is removed, a fresh sweep can collect its report and input. An old publisher which resumes and exposes the report without reacquiring admission then creates a dangling reference. The finite reference records exactly this counterexample. No numerical choice of timeout rules it out without an assumption bounding pauses or a mechanism preventing a stale publisher from committing~\citep{exec17_retention_reference}.

A proposed repair associates the publication right with a generation token. Revocation changes the current generation under coordination; publication must check that its token remains current and that its closure remains present. An old publisher can finish computing, but cannot expose its result using a revoked token. It may start a new admission attempt and reuse still-valid artifacts. This proposal needs a precise failure model before implementation: the generation check and root publication must not themselves be separable by revocation. It is not sufficient to check the token once when a long computation starts.

A local prototype can begin with a simpler stop-the-world collection phase: prevent new admissions, reconcile existing jobs, collect under one owner, then reopen admission. That sacrifices availability but makes the boundary easier to inspect. A concurrent design with publication pins permits more overlap and requires stronger coordination and recovery logic. The choice should depend on the expected workflow and tolerable interruption, not a general preference for concurrency. Neither alternative licenses deletion of records subject to an external retention hold.

\subsection{Retention classes must remain separate from readable resources}
An accepted report root, an unfinished job pin, a reproducibility hold and a custodial hold have different owners and release conditions. The storage design should record which class is retaining each object. A completed report may no longer be exposed to an agent while a review hold still requires its inputs. Conversely, a visible summary does not authorize the agent to read every reachable raw record. Closure determines what must remain available to authorized recovery; it does not grant a new reader permission.

This distinction matters in a deposit analysis that joins a restricted customer extract to an openly distributed method description. Physical deduplication of identical bytes may save space while their custody descriptors remain different. A release condition on one descriptor cannot silently release a hold belonging to another owner. The collector should compute the union of admissible retention obligations, then apply store-specific deletion authority. The proposal cannot specify actual institutional retention periods or assert that indefinite retention is lawful. Those are future owner inputs. Its current contribution is to expose the interfaces at which such inputs change the deletion decision.

Restoring a backup presents the dual problem. Recovered report bytes are insufficient if their descriptors point to absent input or method identities. Before making restored roots visible, compute the dependency closure against the restored inventory, validate the recorded identities and classify missing dependencies. A report with a missing input can be retained as an incomplete historical artifact while remaining ineligible for a reproducibility claim. Recomputing from a newer extract creates a successor analysis, not restoration of the old result. A recovery test should deliberately omit one dependency and verify this distinction.

The resulting prototype protocol is concrete. Freeze an immutable dependency fixture, declare the root and pin lifecycle, enumerate or inject the publication/collection interleavings, interrupt a publisher before pin release, attempt stale-token publication, and restore a partial inventory. Record every admission, rejection, published root and missing dependency. Measure retained bytes and collection delay separately from integrity violations. The finite reference supplies exact expectations for its small core; the added token, custody and backup behaviors still require implementation and their own checks. Such public/generated evidence can establish whether the proposed operating boundary is useful before seeking authentic platform owners and costly banking experiments.
\label{exec17:i4-pins-end}
